\section{When a check is an algorithm}\label{sec:theory}

A rule fitted at the answer is exact at the answer: on every training example, the completed answer $y^{+}$ is a fixed
point of the cold step~\eqref{eq:run}. In logic programming such a fixed point is a supported model, a model of the
rules' Clark completion\cite{clark1978negation}. Run from nothing, a rule computes something else, for rules without
negation its least model, and the two can differ, because at a finished answer facts can support one another in a
circle with nothing to start them. Fages showed that they coincide when the rules are tight: when a level mapping ranks
every fact above the facts that support it\cite{fages1994consistency}. Our rules contain negation (``no better
candidate''), and the matching condition for them is local stratification, under which each negated condition refers
only to facts of lower rank\cite{apt1988towards}. The completion supplies the ranks. What decides the run is whether the
fitted rule uses them, and here a fit at the answer is blind.

\begin{proposition}[the level condition is invisible at the answer]\label{prop:invisible}
Let $c(u,v)$ be a condition that holds on every positive pair of every completed training answer, such as
$\mathrm{depth}(u) < \mathrm{depth}(v)$ for a parent $u$ of $v$. If a rule $R$ is exact at those answers, so is
$R \wedge c$.
\end{proposition}

The proof is one line: $R \wedge c$ accepts every positive, because $c$ does, and rejects every negative that $R$
rejects. So once other conditions reject the same negatives, nothing in the training answers rewards $c$, and a fit that
prefers fewer conditions can drop it. Kept, the condition does the work of the level mapping.

\begin{proposition}[a stratified check produces the answer]\label{prop:stratified}
Let $\ell$ be a level mapping on the items given by the completion, and suppose that the check of every item $v$ reads,
besides the input, only items $u$ with $\ell(u) < \ell(v)$, and is exact at the completed answer. Number the levels
$0,\ldots,L$. Then the cold run from $\notyet$, in which an item is settled once every item below it is, settles level
$\lambda$ by loop $\lambda+1$ and so reproduces the answer within $L+1$ loops, one per level; when the input fixes level
$0$, as it fixes the source, $L$ loops suffice.
\end{proposition}

The proof is by induction on the level: once the items below level $\lambda$ are settled correctly, an item at level
$\lambda$ reads exactly what it reads at the completed answer, where its check is exact. The level condition ``the support lies below'' is what
makes a fitted check of this kind, since it restricts each item's rule to lower levels; keeping every condition that held
on all training positives guarantees that the rule contains it. Our runs evaluate every item at every loop. For breadth-first search, Prim and the
greedy tasks no item is written before its level, because it has no candidate that passes until its support is placed
(Supplementary Note~1). On weighted graphs the run does write provisional parents and replace them, and there the level
condition does a different job.

\begin{proposition}[the weighted cold run is Bellman--Ford's relaxation]\label{prop:bf}
Let the edge weights be positive and the check of a parent be the fitted rule of the shortest-path tasks: $u$ is the
parent of $v$ when $\mathrm{dist}_u < \mathrm{dist}_v$, an edge $u \to v$, $u$ is reached, $v$ is not the source, and no
$k$ has $\mathrm{dist}_k < \mathrm{dist}_v$ and $S(k,v) < S(u,v)$, with $S(u,v) = \mathrm{dist}_u + w(u,v)$ and the
distances recomputed along the pointer each loop. Then at every loop the pointer is a tree rooted at the source, a
reached node stays reached and its distance never rises, and after loop $t$ every distance is at most the length of the
shortest path of at most $t$ edges. With $h$ the largest number of edges on a shortest path, the run holds the true
distances by loop $h$, holds the shortest-path tree by loop $h+1$ (unique when the shortest paths are, as continuous
weights make them almost surely), and comes to rest at most one loop later.
\end{proposition}

The level condition is what keeps the pointer acyclic: a cycle in the new pointer would need a chain of strictly
increasing distances that returns to its start. A reached node $v$ always keeps a passing candidate, since its current
parent $p$ has $\mathrm{dist}_p < \mathrm{dist}_v$ (the weight is positive), and the in-neighbour below $v$ with the
smallest $S$ passes; so every chain of pointers ends at the source, whose distance is $0$. Distances are path sums of the
new tree, so they are never below the true ones, and by induction along the tree each is at most
$\min_k \mathrm{dist}_k + w(k,v)$ of the loop before, Jacobi's form of the Bellman--Ford bound\cite{bellman1958routing}.
Supplementary Note~1 gives the details; Floyd--Warshall runs the same rule once per source.


\begin{proposition}[a finite certificate for every size]\label{prop:finite}
Suppose a rule and the definition of the answer are both conjunctions of a pair's own comparisons and of witness tests
$\exists k : A(u,k) \wedge B(k,v)$, with $W$ witness tests in all, on a class of inputs closed under deleting items.
If the two agree on every input of at most $2 + W$ items, they agree on every input of every size.
\end{proposition}

For a pair $(u,v)$ in any input, delete every item except $u$, $v$ and one witness for each witness test that holds. A
test that holds still holds, and a test that fails still fails, since deleting items never creates a witness; the own
comparisons do not change. So the pair keeps every truth value, and a disagreement at any size would already show on
at most $2 + W$ items. This is the classical preservation of existential conditions under substructures, applied to a fitted table. Count
comparisons are not preserved under deletion, so the certificate covers the rules that use none. For a rule that reads a
completion that is a global closure, such as the depth or the bottleneck, the certificate works on the completed input
$(x, y^{+})$, with the registers given and constrained only by local axioms; it does not recompute the closure on the
smaller input. Section~\ref{sec:proofs} applies it.

On breadth-first search the three parts of the method separate cleanly (Fig.~\ref{fig:theory}). Here we use the
textbook's queue-order search, in which a node's parent is the neighbour that discovered it, and fit on 30 connected
random graphs of 8 to 16 nodes. With the depth as the completion, the fitted rule always names a valid shortest-path
parent, but it names the label's own parent only \nm{bfs.depth.64}\% of the time at 64 nodes. The label breaks ties by
the order of discovery, which depth does not hold. With the place of each node in the order of discovery as the
completion, the rule is exact at the answer. The place is computed from the parent tree alone, and on a partial tree as
well: the source has place 0, a node whose parent is not yet placed has place $\notyet$, and a placed node's place is the
number of placed nodes before it, ordered by their parents' places and then by number. In the cold run a node keeps its
parent once it has one. Run cold from $\notyet$ as fitted, it
names only \nm{bfs.asfitted.64}\% of the parents right, because the fit dropped the condition that the parent is placed
before the child. With that condition kept, it is exact at 64, 256 and 1,024 nodes and settles one level of the search
per loop. The CLRS-30 library's breadth-first search, used in Section~\ref{sec:clrs}, labels each node differently, with
its lowest-numbered neighbour one level up, and there the depth alone is enough. On that task, and on Prim and Dijkstra,
the fitted rule already contains its level condition, and it is exact from nothing with or without the kept set. Keeping every condition is the safeguard for labellings
like the queue order, where the fit can drop it.

A completion of this kind is what distributed computing calls a proof labelling scheme: labels on the nodes that let a
verifier looking only at each node's neighbourhood confirm a global property, as a distance from the root certifies a
breadth-first tree\cite{korman2010proof,goos2016locally}. The fitted check is the verifier, and the cold run computes the
labels. Seen this way, the hints of algorithmic reasoning, which supervise a network on the steps of an algorithm, are
certificates computed from the trace, and completions are certificates computed from the answer. Where the answer
carries no such label, the method needs help. Depth-first search is the clear case. Completed with its discovery and
finish times, its answer is checkable: the fitted check is exact at the finished answer (\nm{dfs.answer}\%). Run cold,
it names \nm{dfs.cold}\% of the parents right, because the state that decides each step of the search, which visited
nodes are still open, is never visible in a finished answer, where every node is closed.
